\chapter{Conclusion}\label{chapter:conclusion}

\section{Summary}\label{sec:conclusion-summary}

This report examined whether evidence-grounded trust can reduce sycophantic collapse in multi-agent large language model debate. The problem is concrete: a confident but wrong majority can pressure a correct minority agent into abandoning the correct answer, and a final majority vote then presents the result as consensus. The proposed response replaces the equal-weight vote with a trust vector that is updated during the debate from claim-level scientific evidence. Each agent's response is decomposed into claims, each claim is routed to a domain-specific source pair, and a verifier assigns supported, contradicted, or unverifiable. Supported and contradicted claims update the agent's weight through logit accumulation, one softmax, and a simplex projection. The final answer is a trust-weighted plurality.

The project produced three artifacts. The first is an implemented pipeline: a three-agent, three-round debate loop; seeded option rendering; an answer parser and scorer; claim decomposition; claim-domain retrieval; a verifier interface; the bounded trust operator; weighted aggregation; injection; and a resumable record writer. The second is a validation harness: a 523-test suite (514 at the P0 release), a 13-check P0 gate, a mock end-to-end path, and the Gate~0 transcript. The third is a pre-registered evaluation plan: four executable arms across three datasets and two seeds, 24 runs, 1,280 GPU-h under the conservative cap, and a six-arm Full target whose placebo and oracle controls are specified under WP-18. These results are engineering evidence; they do not show that trust weighting reduces collapse or improves accuracy.

\section{Limitations}\label{sec:conclusion-limitations}

The most important limitation is verifier quality. The default verifier is lexical, and lexical overlap is not entailment. It can mistake a shared term for support, miss a negation, and misread a numeric claim. The cross-encoder is optional, and the NLI and LLM-judge paths are interfaces rather than validated components. A trust signal built on a weak verifier can amplify a confident wrong majority instead of correcting it. The pre-registration treats verifier accuracy as a moderator for this reason, but the audit that would measure it is not complete.

The second limitation is the evidence layer's coverage. The live PubMed, arXiv, Semantic Scholar, and OpenAlex calls are not exercised by the test suite. Routing and deduplication are tested, but rate limits, source outages, and domain coverage are operational unknowns. A claim with no retrieved passage becomes unverifiable and does not move trust, so sparse domains weaken the mechanism. The report states this as a boundary rather than hiding it.

The third limitation is the experimental scope. Gate~0 covers ten GPQA questions from one dataset, one development stack, and one three-model panel. The matrix is planned for GPQA-Main, MMLU-Pro STEM, and HLE, but it has not run. The injection protocol is a controlled stress test, not natural sycophancy. The debate is fixed at three rounds and three agents, and the model families are three specific open checkpoints. The results, when they exist, will be conditional on those choices.

The fourth limitation is implementation timing. The trust timing is fixed: the debate loop performs a closing pass on the final positions and exposes the final-position trust vector. The runner-side record adapter copies the manifest hash, the final-answer fields, and the final-position trust vector into every record. The six-arm target is documented and pre-registered, but the placebo and oracle arms depend on WP-18. The final stack is pinned by revision and weight total, not by a successful serving run. These are engineering gaps, not evidence of failure, but they bound what the current implementation can support.

\section{Future Work}\label{sec:conclusion-future}

The next step is the final-stack smoke test on the 96GB card, the WP-18 placebo/oracle slice, the 24-item calibration batch, and the verifier human audit. The runner adapter already stamps every record with the manifest hash and the final-answer and trust fields. These steps are prerequisites for a clean run.

The verifier audit is the scientific prerequisite. A set of 150--200 claim-passage pairs, labelled by two annotators, should be used to report balanced accuracy, precision, recall, coverage, and agreement. If the lexical verifier falls below the pre-registered threshold, the cross-encoder or an LLM judge should be substituted and re-audited before the matrix. A matrix run on a verifier that has not passed the audit would produce a number, but not a trustworthy one.

After the audit, the executable matrix can run: four arms, three datasets, two seeds, and the pre-registered analysis. The oracle arm belongs to the six-arm Full target under WP-18; the Core run is the four-arm prefix and does not include it. The primary contrast is B8 versus B8-no-trust. The co-primary endpoint is correct-minority retention. The secondary endpoints are appropriate revision, trust AUROC, trust Brier, verifier accuracy, and citation faithfulness. The six-arm target adds the placebo and oracle controls. A negative or null result is reportable under the pre-registration; it would bound the conditions under which evidence-grounded trust helps, which is itself useful.

The final-stack smoke test, the $N$-scaling and $\alpha/\beta$ ablations, and the human evaluation follow. The human evaluation should use blinded binary judgments rather than a seven-point Likert scale as the primary endpoint. The longer-term goal is a reproducible package with code, configuration, prompts, seeds, dataset identifiers, and the pre-registration, so that another group can repeat the matrix and inspect the evidence behind each answer.

The project has not shown that trust-calibrated deliberation works. It has shown that the mechanism can be built, tested, bounded, and pre-registered, and it has separated the evidence that exists from the evidence that does not. That separation is the condition for a credible test, and the test is the next step.
